\lampiransection{Dokumentasi Pra-Survei Observasi GrabFood}
\label{app:observasi-grabfood}

Lampiran ini menyajikan dokumentasi observasi awal pada halaman GrabFood Rumah Makan Nasi Gerilya. Pengamatan dilakukan pada tanggal 04 April 2026 untuk melihat informasi yang tampak bagi pelanggan, seperti tampilan toko, rating, jumlah penilaian, kisaran harga, menu terpopuler, dan promo yang tersedia. Sebagai pembanding awal, pengamatan juga mencatat dua usaha sejenis, yaitu RM Padang Paripurna dan RM Padang Dendeng Batokok. Informasi pembanding dan kecenderungan opini pelanggan digunakan pada Bab I untuk memperjelas konteks masalah penelitian.

Data pada lampiran ini digunakan sebagai gambaran awal kondisi Nasi Gerilya di GrabFood. Karena bersifat pra-survei, data tersebut belum dimaksudkan sebagai kesimpulan penelitian, melainkan sebagai pendukung latar belakang dan dasar penyusunan instrumen pengumpulan data.

\begin{table}[!htbp]
    \centering
    \caption{Identitas Observasi GrabFood}
    \label{tab:identitas_observasi_grabfood}
    \begin{tabularx}{\textwidth}{|J|Y|}
        \toprule
        \textbf{Aspek}    & \textbf{Keterangan}                                                           \\ \hline
        \midrule
        Nama usaha        & Rumah Makan Nasi Gerilya                                                      \\ \hline
        Platform          & GrabFood                                                                      \\ \hline
        Tanggal observasi & 04 April 2026                                                                 \\ \hline
        Fokus observasi   & Tampilan toko, rating, penilaian, ulasan, harga, menu teratas, promo, dan opini konsumen \\ \hline
        \bottomrule
    \end{tabularx}
    \tablesource{Observasi platform GrabFood pada 04 April 2026; diolah peneliti (2026).}
\end{table}

\begin{table}[!htbp]
    \centering
    \caption{Profil dan Data Utama Rumah Makan Nasi Gerilya}
    \label{tab:data_gerilya}
    \begin{tabularx}{\textwidth}{|J|Y|}
        \toprule
        \textbf{Kategori Data} & \textbf{Fakta / Temuan}                                                                             \\ \hline
        \midrule
        Rating                 & 4,7/5,0                                                                                             \\ \hline
        Total penilaian        & 4.361 penilaian                                                                                     \\ \hline
        Harga terendah         & Rp11.000 (Nasi Putih)                                                                               \\ \hline
        Harga tertinggi        & Rp72.000 (Nasi Padang Kari Kambing Gerilya)                                                         \\ \hline
        Menu teratas           & Nasi Padang Rendang Sapi; Nasi Padang Ayam Goreng; Nasi Padang Ayam Bakar; Nasi Bakar Kikil Tunjang \\ \hline
        Harga menu teratas     & Rp56.000; Rp45.000; Rp45.000; Rp68.000                                                              \\ \hline
        Promo Grab             & Diskon Rp8.000 Jajan Hemat; Diskon 25\% maksimal Rp100.000; Diskon 15\% Group Order                 \\ \hline
        Promo bank             & Raya, BTN, BSI, Mega, Mega Syariah, SBI, OCBC, Danamon, Muamalat                                    \\ \hline
        \bottomrule
    \end{tabularx}
    \tablesource{Observasi platform GrabFood pada 04 April 2026; diolah peneliti (2026).}
\end{table}

Daftar harga disusun dengan menelusuri menu yang tampil pada halaman GrabFood Rumah Makan Nasi Gerilya. Beberapa tangkapan layar dipilih dan ditampilkan pada bagian yang relevan agar informasi utama lebih mudah dibaca, sedangkan ringkasan harga setiap menu disajikan pada Tabel~\ref{tab:survei_harga_nasi_gerilya_grabfood}.

\begin{table}[!htbp]
    \centering
    \caption{Indeks Gambar Pendukung Observasi GrabFood}
    \label{tab:kode-bukti-prasurvei-grabfood}
    \footnotesize
    \setlength{\tabcolsep}{4pt}
    \renewcommand{\arraystretch}{0.95}
    \begin{tabularx}{\textwidth}{|Z{2.4cm}|Z{3.2cm}|Y|}
        \toprule
        \textbf{Kode Gambar} & \textbf{Aspek} & \textbf{Keterangan} \\ \hline
        \midrule
        PA-GF-01 & Rating dan penilaian & Menunjukkan nilai 4,7/5,0 dan jumlah penilaian awal Rumah Makan Nasi Gerilya pada GrabFood. \\ \hline
        PA-HRG-01 & Menu utama & Menunjukkan harga paket nasi Padang dan menu utama yang muncul pada halaman GrabFood. \\ \hline
        PA-HRG-02 & Lauk dan tambahan & Menunjukkan variasi lauk tambahan, pesanan tambahan, dan rentang harga menu pendamping. \\ \hline
        PA-HRG-03 & Sambal dan produk kemasan & Menunjukkan variasi sambal serta produk kemasan yang ditawarkan pada platform. \\ \hline
        PA-HRG-04 & Kuah dan minuman & Menunjukkan item tambahan berharga rendah serta menu minuman yang dapat menambah variasi transaksi. \\ \hline
        PA-HRG-05 & Menu premium & Menunjukkan keberadaan menu premium dengan harga tinggi pada halaman GrabFood. \\ \hline
        \bottomrule
    \end{tabularx}
    \tablesource{Diolah peneliti dari observasi GrabFood tanggal 04 April 2026.}
\end{table}

\clearpage
\subsection*{Dokumentasi Visual Terpilih}

Gambar berikut memperlihatkan bagian utama dari tangkapan layar yang menjadi dasar pencatatan data pada observasi awal. Setiap gambar diberi keterangan dan kode gambar agar pembaca dapat menelusuri hubungan antara gambar dan rekap harga pada Tabel~\ref{tab:survei_harga_nasi_gerilya_grabfood}.

\newlength{\praSurveyEvidenceMaxHeight}
\newcommand{\praSurveyEvidenceBlock}[4]{%
    \setlength{\praSurveyEvidenceMaxHeight}{#4\textheight}%
    \par\noindent\begin{minipage}{\textwidth}%
        \centering
        \includegraphics[width=0.74\textwidth,height=\praSurveyEvidenceMaxHeight,keepaspectratio]{#3}\par
        \figuresource[0.74\textwidth]{Dokumentasi screenshot GrabFood tanggal 04 April 2026; diolah peneliti (2026).}
        \captionof{figure}{#1 -- #2}
    \end{minipage}\par\vspace{8pt}%
}

\praSurveyEvidenceBlock{PA-GF-01}{Rating dan jumlah penilaian awal Rumah Makan Nasi Gerilya pada GrabFood.}{resource/03-prasurvey/grabfood-observation/cropped/PA-GF-01-rating-ulasan.jpeg}{0.60}
\praSurveyEvidenceBlock{PA-HRG-01}{Tampilan menu utama dan harga paket nasi Padang pada halaman GrabFood.}{resource/03-prasurvey/grabfood-observation/cropped/PA-HRG-01-menu-utama.jpeg}{0.60}
\praSurveyEvidenceBlock{PA-HRG-02}{Tampilan variasi lauk dan pesanan tambahan pada halaman GrabFood.}{resource/03-prasurvey/grabfood-observation/cropped/PA-HRG-02-pesanan-tambahan.jpeg}{0.60}
\praSurveyEvidenceBlock{PA-HRG-03}{Tampilan variasi sambal dan produk kemasan pada halaman GrabFood.}{resource/03-prasurvey/grabfood-observation/cropped/PA-HRG-03-sambal-kemasan.jpeg}{0.55}
\praSurveyEvidenceBlock{PA-HRG-04}{Tampilan menu kuah tambahan dan minuman pada halaman GrabFood.}{resource/03-prasurvey/grabfood-observation/cropped/PA-HRG-04-kuah-minuman.jpeg}{0.60}
\praSurveyEvidenceBlock{PA-HRG-05}{Tampilan menu premium kepala ikan pada halaman GrabFood.}{resource/03-prasurvey/grabfood-observation/cropped/PA-HRG-05-menu-premium.jpeg}{0.48}

\clearpage

{\scriptsize
    \setlength{\tabcolsep}{2pt}
    \renewcommand{\arraystretch}{1.08}
    \setlength{\LTcapwidth}{\linewidth}
    \refstepcounter{table}\label{tab:survei_harga_nasi_gerilya_grabfood}
    \edef\SurveiHargaNomor{\thetable}
    \newcommand{\BuktiHargaSatu}{PA-HRG-01}
    \newcommand{\BuktiHargaDua}{PA-HRG-02}
    \newcommand{\BuktiHargaTiga}{PA-HRG-03}
    \newcommand{\BuktiHargaEmpat}{PA-HRG-04}
    \newcommand{\BuktiHargaLima}{PA-HRG-05}
    \addcontentsline{lot}{table}{\protect\numberline{\thetable}Survei Harga Menu Rumah Makan Nasi Gerilya pada GrabFood}
    \begin{longtable}{|>{\raggedright\arraybackslash}p{\dimexpr\textwidth-4.8cm-6\tabcolsep-4\arrayrulewidth\relax}|>{\raggedleft\arraybackslash}p{2.4cm}|p{2.4cm}|}
        \multicolumn{3}{c}{\textbf{Tabel \SurveiHargaNomor}}                                        \\
        \multicolumn{3}{c}{\textbf{Survei Harga Menu Rumah Makan Nasi Gerilya pada GrabFood}}       \\[6pt]
        \toprule
        \textbf{Nama Menu}                                 & \textbf{Harga} & \textbf{Kode Gambar} \\ \hline
        \midrule
        \endfirsthead
        \multicolumn{3}{l}{\textbf{Lanjutan Tabel \SurveiHargaNomor}}                               \\[2pt]
        \toprule
        \textbf{Nama Menu}                                 & \textbf{Harga} & \textbf{Kode Gambar} \\ \hline
        \midrule
        \endhead
        \midrule
        \multicolumn{3}{r}{\footnotesize Dilanjutkan pada halaman berikutnya}                       \\
        \endfoot
        \bottomrule
        \endlastfoot
        \rowcolor{black!8}\multicolumn{3}{|l|}{\textbf{A. GERILYA Nasi Padang}}                     \\ \hline
        Nasi Padang Dendeng Sapi Bakar GERILYA             & Rp56.000       & \BuktiHargaSatu       \\ \hline
        Nasi Padang Kikil Tunjang Gulai GERILYA            & Rp68.000       & \BuktiHargaSatu       \\ \hline
        Nasi Padang Ayam Goreng Bumbu GERILYA              & Rp45.000       & \BuktiHargaSatu       \\ \hline
        Nasi Padang Gulai Ayam Kalasan GERILYA             & Rp46.000       & \BuktiHargaSatu       \\ \hline
        Nasi Padang Ayam Bakar GERILYA                     & Rp45.000       & \BuktiHargaSatu       \\ \hline
        Nasi Padang Udang Goreng Panas GERILYA             & Rp58.000       & \BuktiHargaSatu       \\ \hline
        Nasi Padang Ikan Lele Goreng Panas GERILYA         & Rp39.000       & \BuktiHargaSatu       \\ \hline
        Nasi Padang Ikan Gembung Goreng GERILYA            & Rp50.000       & \BuktiHargaSatu       \\ \hline
        Nasi Padang Ikan Nila Goreng GERILYA               & Rp43.000       & \BuktiHargaSatu       \\ \hline
        Nasi Padang Ikan Nila Bakar Bumbu GERILYA          & Rp50.000       & \BuktiHargaSatu       \\ \hline
        Nasi Padang Ikan Kakap Goreng Panas GERILYA        & Rp48.000       & \BuktiHargaSatu       \\ \hline
        Nasi Padang Telur Gulai GERILYA                    & Rp32.000       & \BuktiHargaSatu       \\ \hline
        Nasi Padang Berkedel Kentang GERILYA               & Rp37.400       & \BuktiHargaSatu       \\ \hline
        Nasi Padang Sayur Polos GERILYA                    & Rp27.500       & \BuktiHargaSatu       \\ \hline
        Nasi Padang Gembung Kuring Bakar GERILYA           & Rp47.850       & \BuktiHargaSatu       \\ \hline
        Nasi Padang Kari Kambing GERILYA                   & Rp72.000       & \BuktiHargaSatu       \\ \hline
        Nasi Padang Ayam Pop Sauce Creamy GERILYA          & Rp46.000       & \BuktiHargaSatu       \\ \hline
        Nasi Padang Rendang Sapi GERILYA                   & Rp56.000       & \BuktiHargaSatu       \\ \hline
        Nasi Padang Ayam Rendang GERILYA                   & Rp46.000       & \BuktiHargaSatu       \\ \hline
        Nasi Padang Cumi Nagih GERILYA                     & Rp70.000       & \BuktiHargaSatu       \\ \hline
        Nasi Padang Ikan Marlin Goreng GERILYA             & Rp51.500       & \BuktiHargaDua        \\ \hline
        Nasi Padang Ikan Tongkol Tuna Goreng Panas GERILYA & Rp48.000       & \BuktiHargaDua        \\ \hline
        Nasi Padang Kakap Muncung Gulai Kuning GERILYA     & Rp51.150       & \BuktiHargaDua        \\ \hline
        Nasi Padang KAKAP MERAH Gulai Kuning GERILYA       & Rp69.300       & \BuktiHargaDua        \\ \hline
        Nasi Padang Telur Dadar GERILYA                    & Rp32.000       & \BuktiHargaDua        \\ \hline
        \rowcolor{black!8}\multicolumn{3}{|l|}{\textbf{B. Pesanan Tambahan}}                        \\ \hline
        Jangek Siram                                       & Rp25.000       & \BuktiHargaDua        \\ \hline
        Keripik Kentang Balado                             & Rp27.500       & \BuktiHargaDua        \\ \hline
        Tambah Lauk Telur Gulai Merah                      & Rp17.500       & \BuktiHargaDua        \\ \hline
        Tambah Lauk Berkedel Kentang (size besar 85g)      & Rp17.500       & \BuktiHargaDua        \\ \hline
        Nasi Putih Tambah                                  & Rp11.000       & \BuktiHargaDua        \\ \hline
        Sate Rendang Jengkol                               & Rp7.700        & \BuktiHargaDua        \\ \hline
        Tambah Lauk Telur Dadar Sayur (1 butir lebih)      & Rp17.500       & \BuktiHargaDua        \\ \hline
        \rowcolor{black!8}\multicolumn{3}{|l|}{\textbf{C. 1 Lauk, Ngak Cukup!}}                     \\ \hline
        Gulai cincang Gerilya                              & Rp55.000       & \BuktiHargaDua        \\ \hline
        Tambah Lauk Ayam Goreng Bumbu                      & Rp35.200       & \BuktiHargaDua        \\ \hline
        Tambah Lauk Gulai Ayam Kalasan                     & Rp36.000       & \BuktiHargaDua        \\ \hline
        Tambah Lauk Ayam Bakar Bumbu                       & Rp35.200       & \BuktiHargaDua        \\ \hline
        Tambah Lauk Gulai Kikil Tunjang                    & Rp63.000       & \BuktiHargaDua        \\ \hline
        Tambah Lauk Dendeng Sapi Bakar                     & Rp43.000       & \BuktiHargaDua        \\ \hline
        Tambah Lauk Udang Goreng Panas                     & Rp46.200       & \BuktiHargaDua        \\ \hline
        Tambah Lauk Ikan Lele Goreng Panas                 & Rp27.000       & \BuktiHargaDua        \\ \hline
        Tambah Lauk Ikan Nila Goreng Bumbu                 & Rp37.400       & \BuktiHargaTiga       \\ \hline
        Tambah Lauk Ikan Kakap Goreng Panas                & Rp38.000       & \BuktiHargaTiga       \\ \hline
        Tambah Lauk Kari Kambing                           & Rp67.500       & \BuktiHargaTiga       \\ \hline
        Tambah Lauk Ikan Gembung Bakar                     & Rp41.250       & \BuktiHargaTiga       \\ \hline
        Tambah Lauk Telur Ikan Sambal Merah GERILYA        & Rp40.150       & \BuktiHargaTiga       \\ \hline
        Tambah Lauk Ayam Pop Creamy GERILYA                & Rp38.500       & \BuktiHargaTiga       \\ \hline
        Tambah Lauk Ikan Marlin Goreng Bumbu               & Rp38.000       & \BuktiHargaTiga       \\ \hline
        Tambah Lauk Ayam Rendang (Kalasan)                 & Rp38.000       & \BuktiHargaTiga       \\ \hline
        Tambah Lauk Rendang Sapi                           & Rp46.000       & \BuktiHargaTiga       \\ \hline
        Tambah Lauk Cumi Nagih                             & Rp63.000       & \BuktiHargaTiga       \\ \hline
        Tambah Lauk Ikan Nila Bakar Bumbu                  & Rp41.000       & \BuktiHargaTiga       \\ \hline
        Tambah Lauk Ikan Gembung Kuring Goreng Panas       & Rp35.000       & \BuktiHargaTiga       \\ \hline
        Tambah Lauk Ikan Tongkol Tuna Goreng Bumbu         & Rp38.000       & \BuktiHargaTiga       \\ \hline
        Tambah Lauk KAKAP MERAH Gulai Kuning GERILYA       & Rp58.300       & \BuktiHargaTiga       \\ \hline
        Tambah Lauk Kakap Muncung Gulai Kuning GERILYA     & Rp40.150       & \BuktiHargaTiga       \\ \hline
        \rowcolor{black!8}\multicolumn{3}{|l|}{\textbf{D. Sambal}}                                  \\ \hline
        GERILYA Sambal Hijau                               & Rp17.500       & \BuktiHargaTiga       \\ \hline
        GERILYA Sambal Merah                               & Rp17.500       & \BuktiHargaTiga       \\ \hline
        \rowcolor{black!8}\multicolumn{3}{|l|}{\textbf{E. Sambal Gerilya Kemasan}}                  \\ \hline
        GERILYA Sambal Andaliman Getir (Kemasan)           & Rp38.500       & \BuktiHargaTiga       \\ \hline
        GERILYA Sambal Dencis Klotok (Kemasan)             & Rp38.500       & \BuktiHargaTiga       \\ \hline
        GERILYA Sambal Teri Medan Kali (Kemasan)           & Rp38.500       & \BuktiHargaEmpat      \\ \hline
        GERILYA Sambal Cumi Nagih (Kemasan)                & Rp38.500       & \BuktiHargaEmpat      \\ \hline
        GERILYA Sambal Tuna Asap Keumamah Aceh (Kemasan)   & Rp38.500       & \BuktiHargaEmpat      \\ \hline
        GERILYA Sambal Bawang Gurih (Kemasan)              & Rp38.500       & \BuktiHargaEmpat      \\ \hline
        \rowcolor{black!8}\multicolumn{3}{|l|}{\textbf{F. Nambah Kuah}}                             \\ \hline
        Nambah KUAH CAMPUR                                 & Rp2.500        & \BuktiHargaEmpat      \\ \hline
        Nambah Kuah Gulai Kuning                           & Rp1.100        & \BuktiHargaEmpat      \\ \hline
        \rowcolor{black!8}\multicolumn{3}{|l|}{\textbf{G. Minuman Segar}}                           \\ \hline
        Air Mineral (400ml)                                & Rp13.200       & \BuktiHargaEmpat      \\ \hline
        Es Teh Tawar                                       & Rp13.200       & \BuktiHargaEmpat      \\ \hline
        Es Teh Manis                                       & Rp15.950       & \BuktiHargaEmpat      \\ \hline
        Es Jeruk Limau                                     & Rp22.000       & \BuktiHargaEmpat      \\ \hline
        Es Lemon Tea                                       & Rp26.400       & \BuktiHargaEmpat      \\ \hline
        Kopi Hitam Dingin                                  & Rp23.100       & \BuktiHargaEmpat      \\ \hline
        Kopi Susu Dingin                                   & Rp26.400       & \BuktiHargaEmpat      \\ \hline
        Jus Jeruk Peras                                    & Rp22.000       & \BuktiHargaEmpat      \\ \hline
        Minuman Marqisa                                    & Rp20.000       & \BuktiHargaEmpat      \\ \hline
        \rowcolor{black!8}\multicolumn{3}{|l|}{\textbf{H. Minuman Panas}}                           \\ \hline
        Teh Tawar Hangat                                   & Rp13.200       & \BuktiHargaEmpat      \\ \hline
        Teh Manis Hangat                                   & Rp15.950       & \BuktiHargaEmpat      \\ \hline
        Lemon Tea Panas                                    & Rp23.100       & \BuktiHargaEmpat      \\ \hline
        Kopi Hitam Panas                                   & Rp23.100       & \BuktiHargaEmpat      \\ \hline
        Kopi Susu Panas                                    & Rp26.400       & \BuktiHargaLima       \\ \hline
        \rowcolor{black!8}\multicolumn{3}{|l|}{\textbf{I. GERILYA Gulai Merah Kepala Ikan}}         \\ \hline
        Gulai Merah Kepala Ikan GERILYA L                  & Rp275.000      & \BuktiHargaLima       \\ \hline
        Gulai Merah Kepala Ikan GERILYA M                  & Rp248.050      & \BuktiHargaLima       \\ \hline
    \end{longtable}
    \tablesource[\textwidth]{Observasi screenshot menu harga Rumah Makan Nasi Gerilya pada GrabFood, 04 April 2026; diolah peneliti (2026).}
    \addtocounter{table}{-1}
}

Dokumentasi pada lampiran ini memperlihatkan kondisi awal Nasi Gerilya di GrabFood, terutama dari sisi tampilan toko, rentang harga, variasi menu, dan promosi yang tersedia. Uraian mengenai opini pelanggan, posisi terhadap kompetitor, dan arah analisis SWOT disampaikan pada bagian utama naskah serta akan diperdalam melalui pengumpulan data penelitian.
